We compare an echo state network (ESN), a one-step MLP on delay coordinates and a small GRU, all reimplemented, for closed-loop forecasting of Lorenz-63 and a 10-variable Lorenz-96 system on a CPU. All three share one data pipeline, validation-based tuning (including the optimiser steps and learning rate of the MLP and GRU), one definition of valid prediction time (VPT, in Lyapunov times) and a Wasserstein-based climate metric of long free runs. With 5000 training samples and 5 seeds the ESN reaches a VPT of \rhval{main/esn/lorenz63/vpt/mean:3} on Lorenz-63 against \rhval{main/mlp-delay/lorenz63/vpt/mean:3} (MLP) and \rhval{main/gru/lorenz63/vpt/mean:3} (GRU), but on Lorenz-96 the MLP (\rhval{main/mlp-delay/lorenz96/vpt/mean:3}) is ahead of the ESN (\rhval{main/esn/lorenz96/vpt/mean:3}; paired $p=\rhval{cmp/main/mlp-delay/lorenz96/vpt/paired_p:3}$), so our hypothesis that the ESN wins on both tasks is refuted. The size of the Lorenz-63 gap depends on the MLP's optimiser budget: the ESN/MLP ratio is \rhval{cmp/st_mlp4000/mlp-delay/lorenz63/vpt/ratio:2} at 4000 Adam steps and \rhval{cmp/st_mlp16000/mlp-delay/lorenz63/vpt/ratio:2} at the tuned 16000 (3 seeds). Across spectral radii 0.1--2, seed-mean ESN VPT ranges from \rhval{sweep_rho/esn@rho=2.0/lorenz63/vpt/mean:2} to \rhval{sweep_rho/esn@rho=0.1/lorenz63/vpt/mean:2} on Lorenz-63 and from \rhval{sweep_rho/esn@rho=2.0/lorenz96/vpt/mean:2} to \rhval{sweep_rho/esn@rho=0.1/lorenz96/vpt/mean:2} on Lorenz-96, and it is higher at reservoir size 1000 than at 100 on both tasks. At 500 training samples the ESN has the highest mean VPT on both tasks, but its margin over the MLP is not significant with 3 seeds. At the main setting the three systems are not separated in attractor statistics; without squared readout features (ESN) or input noise (GRU), almost all long rollouts on Lorenz-96 leave the attractor region. Scale is small (two systems, 3--5 seeds, small tuning grids), so the results are indicative only.
